\section{Correctness}\label{sec:correctness}

When is the path returned by A* a cheapest one? Table~\ref{tab:guarantees}
summarizes the answer, which depends on the heuristic. Throughout this section,
$G = (V,E,w)$ is finite with nonnegative weights, $t$ is reachable from $s$, and
$C^{*} = \delta(s,t)$; the key of a vertex $u$ in $Q$ is $f(u) = g(u) + h(u)$.

\begin{table}[htbp]
  \centering
  \caption{What A* guarantees for each kind of heuristic.}\label{tab:guarantees}
  \renewcommand{\arraystretch}{1.25}%
  \begin{tabular}{@{}P{3.2cm}P{4.1cm}P{3.4cm}P{2.8cm}@{}}
    \toprule
    Heuristic & Algorithm~\ref{alg:astar} (closed set) & A* with reopening
      & Result\\
    \midrule
    consistent & \textcolor{algastar}{shortest path}, each vertex expanded
      at most once & \textcolor{algastar}{shortest path}
      & Theorem~\ref{thm:consistent}\\
    admissible, not consistent & \textcolor{cbRed}{may fail}
      (Example~\ref{ex:inconsistent}) & \textcolor{algastar}{shortest path}
      & Theorem~\ref{thm:admissible}\\
    not admissible & \textcolor{cbRed}{may fail}
      (Section~\ref{sec:overestimate}) & \textcolor{cbRed}{may fail} & none\\
    $h(v) = 0$ for all $v$ (Dijkstra's algorithm)
      & \textcolor{algastar}{shortest path} & \textcolor{algastar}{shortest path}
      & Corollary~\ref{cor:dijkstra}\\
    \bottomrule
  \end{tabular}
\end{table}

\subsection{Two facts about the cost so far}\label{sec:invariant}

The values $g(v)$ only go down and never drop below the true distance from the
source. The same proof applies to A* with reopening
(Section~\ref{sec:admissible-proof}). A \emph{walk} is like a path but may
visit a vertex more than once.

\begin{lemma}\label{lem:invariant}
During the execution of Algorithm~\ref{alg:astar}, the value $g(v)$ of a vertex
$v$ never increases. Whenever $g(v)$ is finite, it is the cost of some walk from
$s$ to $v$, and consequently $g(v) \ge \delta(s,v)$.
\end{lemma}

\begin{proof}
The algorithm changes $g(v)$ only by the assignment $g(v) \gets g(u) + w(u,v)$,
and only when the new value is smaller, so $g(v)$ never increases. We prove the
second claim by induction on the number of assignments. Initially only $g(s)$ is
finite, and $g(s) = 0$ is the cost of the walk $(s)$. An assignment sets
$g(v) = g(u) + w(u,v)$, where, by the induction hypothesis, $g(u)$ is the cost
of a walk from $s$ to $u$. Appending $v$ to this walk gives a walk from $s$ to
$v$ of cost $g(v)$. Finally, removing the repeated parts of a walk yields a path
between the same vertices, and the path costs no more than the walk because
every weight is nonnegative. Hence $g(v) \ge \delta(s,v)$.
\end{proof}

\subsection{Consistent heuristics}\label{sec:consistent-proof}

With a consistent heuristic, the estimated total cost never decreases along a
path.

\begin{lemma}\label{lem:pathmono}
Let $h$ be a consistent heuristic and let $P = (v_0, v_1, \dots, v_k)$ be a
path. For $0 \le i \le k$, let $f_P(v_i) = c(v_0, \dots, v_i) + h(v_i)$, where
$c(v_0, \dots, v_i)$ is the cost of the first $i$ edges of $P$. Then
$f_P(v_0) \le f_P(v_1) \le \dots \le f_P(v_k)$.
\end{lemma}

\begin{proof}
For $1 \le i \le k$, Inequality~\eqref{eq:consistent} applied to the edge
$(v_{i-1}, v_i)$ gives
\begin{equation}\label{eq:fstep}
  f_P(v_i) = f_P(v_{i-1}) + w(v_{i-1}, v_i) + h(v_i) - h(v_{i-1})
           \ge f_P(v_{i-1}). \qedhere
\end{equation}
\end{proof}

Consequently, the extracted keys never decrease.

\begin{lemma}\label{lem:monotone}
Let $h$ be a consistent heuristic. Then Algorithm~\ref{alg:astar} extracts
vertices from $Q$ in nondecreasing order of their keys.
\end{lemma}

\begin{proof}
Suppose that the algorithm extracts a vertex $u$ with key $f(u)$. At that moment
every key in $Q$ is at least $f(u)$, because \textsc{Extract-Min} returns a
smallest key. While $u$ is expanded, a key is inserted or decreased only to a
value $g(u) + w(u,v) + h(v)$ for a neighbor $v$ of $u$, and by
Inequality~\eqref{eq:consistent} this value is at least $g(u) + h(u) = f(u)$.
Hence every key in $Q$ is still at least $f(u)$ at the next extraction, and the
next extracted key is at least $f(u)$.
\end{proof}

The main lemma states that a vertex leaves $Q$ only when its value $g$ is exact
(Figure~\ref{fig:proof-idea}).

\begin{figure}[htbp]
  \centering
  \begin{tikzpicture}[x=1cm, y=1cm]
  \fill[gray!12, rounded corners=4pt] (-0.5,-0.55) rectangle (4.5,0.55);
  \node[font=\scriptsize, text=inkgray, anchor=south] at (2,0.55) {closed vertices (in $S$)};
  \node[storystart] (s) at (0,0) {$s$};
  \node[storynode, fill=gray!30, draw=gray] (a) at (2,0) {};
  \node[storynode, fill=gray!30, draw=gray] (y) at (4,0) {$y$};
  \node[storynode, fill=rbluelt] (x) at (6,0) {$x$};
  \node[storynode, fill=hlyellow, draw=alggreedy] (u) at (10,0) {$u$};
  \draw[storyedge, line width=1.2pt] (s) -- (a) -- (y) -- (x);
  \draw[storyedge, line width=1.2pt, dashed] (x) -- node[above, font=\scriptsize,
    text=inkgray] {rest of $P$} (u);
  \node[font=\scriptsize, text=rbluedk, below=10pt, align=center] at (x)
    {first vertex not in $S$\\in $Q$ with $g(x) = \delta(s,x)$};
  \node[font=\scriptsize, text=alggreedy!80!black, below=10pt, align=center] at (u)
    {being extracted\\supposed $g(u) > \delta(s,u)$};
  \node[draw=cbRed, fill=cbRed!5, font=\small, rounded corners=2pt, inner sep=4pt]
    at (8,1.25) {$f(x) \le \delta(s,u) + h(u) < f(u)$};
\end{tikzpicture}

  \caption{The idea behind Lemma~\ref{lem:exact}. On a shortest path $P$ from
    $s$ to $u$, the first vertex $x$ that is not closed is in $Q$ with an exact
    value $g(x)$, and its key is smaller than the key of $u$.}\label{fig:proof-idea}
\end{figure}

\begin{lemma}\label{lem:exact}
Let $h$ be a consistent heuristic. Whenever Algorithm~\ref{alg:astar} extracts a
vertex $u$ from $Q$, we have $g(u) = \delta(s,u)$.
\end{lemma}

\begin{proof}
Suppose the contrary, and let $u$ be the first extracted vertex with
$g(u) \ne \delta(s,u)$; by Lemma~\ref{lem:invariant}, $g(u) > \delta(s,u)$. The
first iteration extracts $s$ with $g(s) = 0 = \delta(s,s)$, so $u \neq s$. The
value $g(u)$ is finite, so $u$ is reachable from $s$; let $P$ be a shortest path
from $s$ to $u$. Just before $u$ is extracted, $s$ is in $S$ and $u$ is not. Let
$x$ be the first vertex of $P$ that is not in $S$, and let $y$ be the vertex
that precedes $x$ on $P$; then $y$ is in $S$.

The vertex $y$ was extracted before $u$, so $g(y) = \delta(s,y)$ by the choice
of $u$. When $y$ was expanded, the vertex $x$ was not in $S$, because $S$ only
grows, and the relaxation of the edge $(y,x)$ ensured that
$g(x) \le \delta(s,y) + w(y,x) = \delta(s,x)$; the equality holds because the
part of $P$ from $s$ to $x$ is a shortest path. By Lemma~\ref{lem:invariant},
$g(x)$ has not increased since and is at least $\delta(s,x)$, so
$g(x) = \delta(s,x)$. A vertex with a finite value $g$ that is not in $S$ is in
$Q$, so $x$ is in $Q$. Moreover $x \neq u$, because $g(u) > \delta(s,u)$.

Lemma~\ref{lem:pathmono}, applied to $P$, compares the keys of $x$ and $u$. The
costs of the parts of $P$ that end at $x$ and at $u$ are $\delta(s,x)$ and
$\delta(s,u)$, so
\begin{equation}\label{eq:xbeforeu}
  f(x) = \delta(s,x) + h(x) \le \delta(s,u) + h(u) < g(u) + h(u) = f(u).
\end{equation}
The vertex $x$ is in $Q$ with a smaller key than $u$, which contradicts the
choice of $u$ by \textsc{Extract-Min}.
\end{proof}

We can now prove the main result.

\begin{theorem}[Optimality of A* with a consistent heuristic]\label{thm:consistent}
Let $G$ be a finite graph with nonnegative edge weights, let $s$ and $t$ be
vertices of $G$ such that $t$ is reachable from $s$, and let $h$ be a
consistent heuristic. Then Algorithm~\ref{alg:astar} expands every vertex at
most once and returns a path from $s$ to $t$ of cost $\delta(s,t)$.
\end{theorem}

\begin{proof}
A vertex enters $S$ when it is extracted and never leaves $S$, and the loop over
the edges never inserts a vertex of $S$ into $Q$. Hence every vertex is extracted
at most once, and the algorithm stops after at most $|V|$ iterations.

The algorithm does not return failure. Suppose, to the contrary, that $Q$
becomes empty before $t$ is extracted. Consider a path from $s$ to $t$; its
first vertex $s$ is in $S$ and its last vertex $t$ is not. Let $x$ be the first
vertex of the path that is not in $S$. The vertex that precedes $x$ on the path
is in $S$, and its expansion gave $x$ a finite value $g(x)$ and inserted $x$
into $Q$. The vertex $x$ can leave $Q$ only by being extracted, after which it
would be in $S$ or the algorithm would have returned. Both cases contradict the
assumptions, so the algorithm extracts $t$, and $g(t) = \delta(s,t)$ at that
moment by Lemma~\ref{lem:exact}.

It remains to show that the returned path has cost $g(t)$. Let $v \neq s$ be a
vertex on the path and let $p = \pi(v)$. The last change of $g(v)$ and $\pi(v)$
happened during the expansion of $p$ and set $g(v) = g(p) + w(p,v)$. The value
$g(p)$ did not change afterwards, because $p$ was in $S$. The vertex $p$ was
extracted before $v$, so following the predecessors from $t$ visits vertices in
strictly decreasing order of extraction and therefore reaches $s$. Adding the
equalities $g(v) - g(\pi(v)) = w(\pi(v), v)$ along the path shows that its cost
is $g(t) - g(s) = \delta(s,t)$.
\end{proof}

Dijkstra's algorithm is the special case of the zero heuristic.

\begin{corollary}\label{cor:dijkstra}
Let $G$ be a finite graph with nonnegative edge weights, and let $s$ and $t$ be
vertices of $G$ such that $t$ is reachable from $s$. Then
Algorithm~\ref{alg:dijkstra} returns a shortest path from $s$ to $t$.
\end{corollary}

\begin{proof}
The zero heuristic, $h(v) = 0$ for every vertex $v$, is consistent, because
$0 \le w(u,v) + 0$ for every edge $(u,v)$. With this heuristic the key $g + h$
equals $g$, so Algorithm~\ref{alg:astar} is Algorithm~\ref{alg:dijkstra}. The
claim follows from Theorem~\ref{thm:consistent}.
\end{proof}

\subsection{Admissible heuristics}\label{sec:admissible-proof}

Admissibility alone does not make Algorithm~\ref{alg:astar} correct.

\begin{example}\label{ex:inconsistent}
Consider the graph of Figure~\ref{fig:counterexample} with source $s$ and destination
$t$. The shortest path is $(s, a, c, t)$, of cost 5. The true remaining costs
are $h^{*}(s) = 5$, $h^{*}(a) = 4$, $h^{*}(b) = 5$ and $h^{*}(c) = 3$, so the
heuristic shown in the figure is admissible. It is not consistent, because
$h(a) = 4 > w(a,c) + h(c) = 2$. Table~\ref{tab:counterexample} traces
Algorithm~\ref{alg:astar} on this graph. In step 4 the edge $(a,c)$ offers $c$
the smaller value $1 + 1 = 2$, but $c$ is already closed, so the algorithm skips
the edge and returns $(s, b, c, t)$, of cost 6.
\end{example}

\begin{figure}[htbp]
  \centering
  \begin{tikzpicture}[x=1.35cm, y=1.1cm]
  \coordinate (s) at (0,0);
  \coordinate (a) at (1.8,1.0);
  \coordinate (b) at (1.8,-1.0);
  \coordinate (c) at (3.6,0);
  \coordinate (t) at (5.4,0);
  \draw[storyedge] (s) -- node[storyweight] {1} (a);
  \draw[storyedge] (a) -- node[storyweight] {1} (c);
  \draw[storyedge] (s) -- node[storyweight] {1} (b);
  \draw[storyedge] (b) -- node[storyweight] {2} (c);
  \draw[storyedge] (c) -- node[storyweight] {3} (t);
  \node[storystart] at (s) {$s$};
  \node[storyh, above=11pt] at (s) {$h=0$};
  \node[storynode] at (a) {$a$};
  \node[storyh, above=11pt] at (a) {$h=4$};
  \node[storynode] at (b) {$b$};
  \node[storyh, below=11pt] at (b) {$h=1$};
  \node[storynode] at (c) {$c$};
  \node[storyh, above=11pt] at (c) {$h=1$};
  \node[storygoal] at (t) {$t$};
  \node[storyh, above=11pt] at (t) {$h=0$};
\end{tikzpicture}

  \caption{An admissible heuristic that is not consistent. Edge weights are on
    the edges, and heuristic values are beside the vertices.}\label{fig:counterexample}
\end{figure}

\begin{table}[htbp]
  \centering
  \caption{Algorithm~\ref{alg:astar} on the graph of
    Figure~\ref{fig:counterexample}.}\label{tab:counterexample}
  \begin{tabular}{@{}clll@{}}
    \toprule
    Step & Extracted (key) & Updates & State of $c$\\
    \midrule
    1 & $s$ (0) & $g(a) = 1$, key 5; \ $g(b) = 1$, key 2 & not reached\\
    2 & $b$ (2) & $g(c) = 3$, key 4 & in $Q$, $g(c) = 3$\\
    3 & $c$ (4) & $g(t) = 6$, key 6 & closed, $g(c) = 3$\\
    4 & $a$ (5) & \textcolor{cbRed}{edge $(a,c)$ skipped, $c$ is closed}
      & closed, $g(c) = 3$\\
    5 & $t$ (6) & returns $(s,b,c,t)$, cost \textcolor{cbRed}{6} & \\
    \bottomrule
  \end{tabular}
\end{table}

Here $c$ was closed before $g(c)$ was exact, and the closed set blocked the
correction. \emph{A* with reopening} removes the condition ``$v$ is not in $S$''
from Algorithm~\ref{alg:astar} and moves every vertex whose $g$ decreases back
from $S$ into $Q$. On Example~\ref{ex:inconsistent} it reopens $c$ with key 3 and
returns $(s,a,c,t)$ of cost 5.

\begin{theorem}[Optimality of A* with an admissible heuristic]\label{thm:admissible}
Let $G$ be a finite graph in which every edge weight is positive, let $s$ and $t$
be vertices of $G$ such that $t$ is reachable from $s$, and let $h$ be an
admissible heuristic. Then A* with reopening terminates and returns a path from
$s$ to $t$ of cost $\delta(s,t)$.
\end{theorem}

\begin{proof}
We first prove termination. Let $w_{\min} > 0$ be the smallest edge weight, let
$v$ be a vertex, and let $g_{0}$ be the first finite value of $g(v)$. By
Lemma~\ref{lem:invariant}, every later value of $g(v)$ is the cost of a walk
from $s$ to $v$ of cost at most $g_{0}$. Such a walk has at most
$g_{0} / w_{\min}$ edges, so there are finitely many such walks, and $g(v)$
takes finitely many values. After an extraction, the vertex $v$ returns to $Q$
only if $g(v)$ decreases, so $v$ is extracted at most once for each value of
$g(v)$. Hence the number of iterations is finite.

Let $P^{*} = (s = v_0, v_1, \dots, v_k = t)$ be a shortest path from $s$ to
$t$. We show that before every extraction that precedes the extraction of $t$,
some vertex $v_j$ of $P^{*}$ is in $Q$ with $g(v_j) = \delta(s,v_j)$. In
particular $Q$ is not empty, so the algorithm does not return failure. Let $j$
be the largest index with $g(v_j) = \delta(s,v_j)$; it exists because
$g(s) = 0$. Suppose that $v_j$ is not in $Q$. Since $g(v_j)$ is finite, $v_j$
has been reached, so it is in $S$; in particular $v_j \neq t$, because the
extraction of $t$ ends the search. The value $g(v_j)$ did not decrease after the
last expansion of $v_j$, because a decrease would have reopened $v_j$. Hence
$g(v_j) = \delta(s,v_j)$ during that expansion, which relaxed the edge
$(v_j, v_{j+1})$ and set
$g(v_{j+1}) \le \delta(s,v_j) + w(v_j,v_{j+1}) = \delta(s,v_{j+1})$. By
Lemma~\ref{lem:invariant}, $g(v_{j+1}) = \delta(s,v_{j+1})$ from then on, which
contradicts the choice of $j$. So $v_j$ is in $Q$.

When $t$ is extracted, its key is $g(t) + h(t) = g(t)$, because admissibility
gives $0 \le h(t) \le h^{*}(t) = 0$. Just before this extraction, the vertex
$v_j$ found above is in $Q$ with key
\begin{equation}\label{eq:admbound}
  \delta(s,v_j) + h(v_j) \le \delta(s,v_j) + h^{*}(v_j) = C^{*},
\end{equation}
where the equality holds because $v_j$ lies on a shortest path from $s$ to $t$.
\textsc{Extract-Min} chose $t$, so $g(t) \le C^{*}$. Together with
$g(t) \ge C^{*}$ from Lemma~\ref{lem:invariant}, this gives $g(t) = C^{*}$.

Finally, consider the returned path. For every vertex $v \neq s$ on it, the last
change of $g(v)$ and $\pi(v)$ set $g(v) = g(\pi(v)) + w(\pi(v), v)$, and
$g(\pi(v))$ can only have decreased since. Hence
$g(\pi(v)) + w(\pi(v), v) \le g(v)$. Every weight is positive, so $g$ strictly
decreases along the predecessors; the path therefore visits no vertex twice and
ends at $s$. Adding the inequalities along the path shows that its cost is at
most $g(t) - g(s) = C^{*}$. No path costs less than $C^{*}$, so the cost of the
returned path is exactly $C^{*}$.
\end{proof}

Theorem~\ref{thm:admissible} is the original result of Hart, Nilsson and
Raphael~\cite{hart1968formal}; Theorem~\ref{thm:consistent} covers the usual
implementation~\cite{russell2020aima}, in which each vertex is expanded at most
once. With an inconsistent heuristic a vertex may be expanded many times
(Section~\ref{sec:complexity}).
